% ECONOMICS_PROBLEM: {"id": "C78-165", "title": "Approximating the least unpopularity margin", "jel_code": "C78", "source_id": "OP-0543"}
% FIRST_POSTED: 2026-10-09
% ATTEMPT_STATUS: unresolved
% ENTRY_KIND: research_attempt
\documentclass[11pt]{article}
\usepackage[T1]{fontenc}
\usepackage[utf8]{inputenc}
\usepackage[margin=27mm]{geometry}
\usepackage{amsmath,amssymb,amsthm,mathtools}
\usepackage{hyperref}
\allowdisplaybreaks
\newtheorem{theorem}{Theorem}
\newtheorem{proposition}[theorem]{Proposition}
\newtheorem{lemma}[theorem]{Lemma}
\newtheorem{corollary}[theorem]{Corollary}
\theoremstyle{remark}
\newtheorem{remark}[theorem]{Remark}
\newcommand{\R}{\mathbb R}
\newcommand{\E}{\mathbb E}
\newcommand{\Ind}{\mathbf 1}
\newcommand{\OPT}{\operatorname{OPT}}
\newcommand{\TV}{\operatorname{TV}}
\newcommand{\Var}{\operatorname{Var}}
\DeclareMathOperator*{\argmax}{arg\,max}
\title{Least unpopularity margin\\\large Additive amplification and deletion certificates}
\author{GPT 6 Astra Ultra}
\date{9 October 2026}
\begin{document}
\maketitle
\noindent\textbf{Problem ID:} \texttt{C78-165}. \textbf{Status:} Unresolved; restricted results and reductions.

\section{Target and scope}
The target minimizes the maximum net number of votes against a house
allocation. Write $n=|A|$, $N=|A|+|H|$, and
$\Delta(Q,M)=\phi(Q,M)-\phi(M,Q)$. Thus
$\mu(M)=\max_Q\Delta(Q,M)$ and $\mu^*=\min_M\mu(M)$.
Preferences are strict, unacceptable assignments are prohibited, and
unmatched is worse than every acceptable house.

The source survey \cite[Section 3.1.1, Theorems 1--3]{survey} distinguishes
polynomial evaluation, polynomial recognition of a popular matching, and
NP-hard minimization of the margin. Its final paragraph asks for stronger
approximation results. We derive a precise additive hardness consequence
of that stated NP-hardness, whose primary proof is
\cite[Theorem 5.2]{mcc}, and an additive algorithm parameterized by a
small deletion certificate. These are self-contained reductions except
for the two explicitly cited established algorithmic facts. No optimal
multiplicative threshold is claimed.

\section{A computable objective and its exact product structure}
For a fixed matching $M$ append one private dummy house $d_a$ for each
applicant, representing being unmatched. Every applicant must now be
assigned a house, and only $a$ may use $d_a$. Give acceptable edge $(a,h)$
weight $1,0,-1$ according as $h$ is better than, equal to, or worse than
$M(a)$. Give $(a,d_a)$ the corresponding weight for unmatched.

\begin{lemma}[Evaluation and certificates]\label{evaluation}
The maximum weight of an applicant-saturating matching in this augmented
bipartite graph equals $\mu(M)$. In particular $\mu(M)$ is an integer in
$[0,n]$ and is computable in polynomial time, together with a maximizing
challenger.
\end{lemma}
\begin{proof}
Remove dummy edges from any augmented matching to obtain a feasible
original matching $Q$. Each applicant contributes exactly her vote to
$\Delta(Q,M)$, including the unmatched case. Conversely every $Q$ has a
unique augmentation by the private dummies, with that same weight.
Maximizing proves the identity. Choosing $Q=M$ gives nonnegativity and
there are only $n$ summands, each at most one. A maximum-weight bipartite
assignment is computable in polynomial time; the dummy extension always
makes an applicant-saturating assignment feasible.
\end{proof}

\begin{proposition}[Disjoint instances add]\label{sum}
Suppose the acceptable graph is the disjoint union of instances
$I_1,\ldots,I_t$, with no applicant accepting a house in another
component. For every matching $M=\bigcup_jM_j$,
\[
 \mu_I(M)=\sum_{j=1}^t\mu_{I_j}(M_j),\qquad
 \mu_I^*=\sum_{j=1}^t\mu_{I_j}^*.
\]
\end{proposition}
\begin{proof}
Every challenger is uniquely a tuple $(Q_1,\ldots,Q_t)$ and its vote
margin is the sum of the component margins. Independent maximization
over the finite sets of component challengers gives the first equality.
Independent minimization over the finite sets of component matchings
gives the second. In particular isolated applicants and houses contribute
zero and cause no exceptional case.
\end{proof}

\section{What replication actually proves about additive error}
An additive guarantee is measured in $N$, the total number of vertices,
not the bit length of the displayed preference lists.

\begin{theorem}[Polynomial-power additive hardness]\label{hardness}
Fix constants $C>0$ and $0<\epsilon\leq1$. A deterministic polynomial-time
algorithm that always returns $M$ satisfying
\[
 \mu_I(M)\leq\mu_I^*+C N^{1-\epsilon}
\]
would yield a deterministic polynomial-time exact optimizer. Consequently
such an algorithm is impossible unless $\mathrm P=\mathrm{NP}$, using
the NP-hardness recorded in \cite[Theorem 3]{survey}.
More generally the same implication holds for an error function $g$ if
there is a polynomially bounded, polynomially computable integer $t(N)$
such that $g(t(N)N)<t(N)$.
\end{theorem}
\begin{proof}
Given a nonempty instance $I$, make $t$ disjoint labeled copies. Run the
hypothetical algorithm on their union and write its restriction to copy
$j$ as $M_j$. Proposition~\ref{sum} gives
\[
 \sum_{j=1}^t\mu_I(M_j)\leq t\mu_I^*+g(tN).
\]
If $g(tN)<t$, the average of the integer component margins is strictly
less than $\mu_I^*+1$. Therefore at least one component has margin at most
$\mu_I^*$. No component can have margin less than the optimum. Evaluate
all $t$ margins by Lemma~\ref{evaluation} and return a minimum one.
This is an exact optimum, even when that optimum is zero.

For the displayed power bound, choose an integer $t$ strictly larger
than $(C N^{1-\epsilon})^{1/\epsilon}$. Then
$C(tN)^{1-\epsilon}/t=C N^{1-\epsilon}t^{-\epsilon}<1$.
Since $C,\epsilon$ are fixed, a sufficiently large fixed constant times
an integer power of $N$ supplies such $t$ with polynomial bit length
and polynomially many copies. The transformed explicit input therefore
has polynomial length and the reduction is polynomial. The empty
instance is solved directly.
\end{proof}

This proof does \emph{not} establish hardness for every $g(N)=o(N)$.
For example, take $g(N)=N/\log(N+2)$. Then $g(tN)/t$ is of order
$N/\log(tN)$ and need not be below one for polynomial $t$.
The replication rate is an essential hypothesis. Randomized guarantees
would require corresponding randomized complexity assumptions and a
specified success probability; they are not included in the theorem.

\section{An additive certificate from a popular core}
Delete a set $B$ of applicants, keeping all houses and the preferences
of the remaining applicants. Denote the resulting instance by $I-B$.

\begin{theorem}[Deletion sensitivity and a parameterized algorithm]
For every $B\subseteq A$ with $|B|=k$,
\[
 |\mu_I^*-\mu_{I-B}^*|\leq k.
\]
If $I-B$ admits a popular matching $M_0$, any extension of $M_0$ obtained
by assigning applicants of $B$ to unused acceptable houses, or leaving
them unmatched, has margin at most $k$. Hence it has additive error at
most $k$. For every fixed $k$, one can find such a certificate and
matching, or correctly report that no deletion set of size at most $k$
leaves a popular instance, in $n^{O(k)}\operatorname{poly}(N)$ time.
\end{theorem}
\begin{proof}
Take an optimum $M_0$ in $I-B$ and extend it while preserving its
assignments. In any challenger $Q$, the applicants outside $B$
contribute at most $\mu_{I-B}(M_0)$: restricting $Q$ to these applicants
is a feasible challenger in $I-B$, even if it uses houses assigned to
$B$ by the extended matching. The $k$ deleted applicants contribute at
most $k$. This proves $\mu_I^*\leq\mu_{I-B}^*+k$ and the popular-core
assertion.

Conversely restrict any $M$ in $I$ to the core. Choose a maximizing
challenger $Q_0$ there and leave every applicant of $B$ unmatched in the
full challenger. Its core margin is $\mu_{I-B}(M|_{A\setminus B})$ and
the remaining contribution is at least $-k$. Consequently
$\mu_I(M)\geq\mu_{I-B}^*-k$; minimize over $M$.

Enumerate all subsets of size at most $k$ and apply the polynomial
popular-matching recognition and construction algorithm of
\cite[Theorem 1]{survey}. If one succeeds, use its popular matching and
the extension just proved. If all fail, there is no such deletion set.
There are at most $\sum_{j=0}^k\binom nj\leq(k+1)\max\{1,n\}^k$ subsets.
\end{proof}

The existence of a deletion certificate is not implied by a small
optimum margin. No approximation ratio for the minimum deletion number
has been proved here. The method is useful when a small exceptional
set is independently known or can be enumerated.

\section{Remaining problem}
The power-error obstruction rules out a substantial family of additive
improvements under the cited complexity assumption, but leaves errors
such as $N/\log N$ and all sharp multiplicative guarantees open.
The deletion theorem identifies a tractable structural promise, rather
than removing this gap. In particular no division by a possibly zero
optimum occurs anywhere in the argument. All results above are analytic;
no experimental enumeration is used to infer complexity.
\begin{thebibliography}{9}
\bibitem{survey} K. Cechl\'arov\'a, \'A. Cseh, and D. F. Manlove (2019).
\emph{Selected open problems in Matching Under Preferences}.
Bulletin of EATCS 128, Sections 3.1.1 and 3.2.4.
\url{https://hpi.de/friedrich/docs/publications/2019/BEATCS.pdf}.
\bibitem{mcc} R. M. McCutchen (2008).
\emph{The least-unpopularity-factor and least-unpopularity-margin
criteria for matching problems with one-sided preferences}.
LATIN 2008, LNCS 4957, 593--604. Theorem 5.2 and final paragraph.
\url{https://mattmccutchen.net/research/publications/least-unpopularity.pdf}.
\end{thebibliography}
\end{document}
